
\documentclass[addpoints]{exam}
\usepackage{amsmath}
\usepackage{amssymb}
\usepackage{color}

% \printanswers

\newcommand{\event}{Applications of Derivatives}
\newcommand{\tournament}{}
\def\type{0}

\setlength\fillinlinelength{\textwidth}
\CorrectChoiceEmphasis{\color{red}\bfseries}
\SolutionEmphasis{\color{red}\bfseries}

\setlength\answerclearance{2pt}
\pagestyle{headandfoot}
\runningheadrule
\runningheader{\event}{\tournament}{\ifnum\type=1 Class Set Test \else Full Name:\hspace{1cm} \fi}
\runningfootrule
\runningfooter{}{Page \thepage\ of \numpages}{}

\begin{document}
\begin{center}
    {\huge Grade 12 Calculus \& Vectors \\ \event
    \ifprintanswers \space Key
    \else \space \ifnum\type<2 Test \fi \ifnum\type=1 \textbf{(CLASS SET)} \fi
          \ifnum\type=2 Answer Sheet \fi
    \fi \\ \vspace{3mm}\tournament\vspace{1mm}}
\end{center}

\vspace{.25\textheight}

\begin{center}
    First Name: \ifprintanswers
    \underline{\hspace{3.5cm}}\textbf{KEY}\underline{\hspace{5.5cm}}\\\vspace{5mm}
    \else \underline{\hspace{10cm}}\\ \vspace{5mm} \fi
    Last Name: \underline{\hspace{10cm}}\\ \vspace{5mm}
\end{center}

\noindent\textbf{\underline{Directions:}}
\begin{itemize}
    \ifnum\type=1 \item \textbf{This test is a class set.} Do not write on this paper. \fi
    \item Show all work. Justify all conclusions with calculus (sign charts or derivative tests).
    \item The test is designed to be completed in 75 minutes.
\end{itemize}
\begin{center}
\textit{For grading use only}\\
\multirowgradetable{1}[pages]
\end{center}

\newpage
\begin{questions}

\section*{Part A --- Multiple Choice (15 marks)}
\vspace{5mm}

\question[1] A critical point of $f$ occurs where:
\begin{choices}
    \choice $f(x) = 0$
    \correctchoice $f'(x) = 0$ or $f'(x)$ is undefined
    \choice $f''(x) = 0$
    \choice $f(x)$ reaches its maximum value
\end{choices}

\question[1] If $f'(x) > 0$ on an interval, then $f$ is:
\begin{choices}
    \choice Concave up
    \choice Decreasing
    \correctchoice Increasing
    \choice At a maximum
\end{choices}

\question[1] If $f''(x) > 0$ on an interval, the curve is:
\begin{choices}
    \choice Increasing
    \choice Decreasing
    \correctchoice Concave up
    \choice Concave down
\end{choices}

\question[1] An inflection point of $f$ occurs where:
\begin{choices}
    \choice $f'(x) = 0$
    \correctchoice $f''(x) = 0$ and $f''$ changes sign
    \choice $f(x) = 0$
    \choice $f''(x) > 0$
\end{choices}

\question[1] If $f'(x)$ changes from positive to negative at $x = c$, then $f$ has a:
\begin{choices}
    \choice Local minimum at $x = c$
    \choice Inflection point at $x = c$
    \correctchoice Local maximum at $x = c$
    \choice Vertical asymptote at $x = c$
\end{choices}

\question[1] The second derivative test: if $f'(c) = 0$ and $f''(c) > 0$, then $x = c$ is a:
\begin{choices}
    \choice Local maximum
    \choice Inflection point
    \correctchoice Local minimum
    \choice The test is inconclusive
\end{choices}

\question[1] A farmer has 100 m of fencing to enclose a rectangular pen. What dimensions
maximize the area?
\begin{choices}
    \choice $20 \text{ m} \times 30 \text{ m}$
    \correctchoice $25 \text{ m} \times 25 \text{ m}$
    \choice $10 \text{ m} \times 40 \text{ m}$
    \choice $15 \text{ m} \times 35 \text{ m}$
\end{choices}

\question[1] The absolute maximum of $f(x) = x^3 - 3x$ on $[-2,\, 2]$ is:
\begin{choices}
    \choice $-2$
    \choice $0$
    \correctchoice $2$
    \choice $4$
\end{choices}

\question[1] The Mean Value Theorem guarantees the existence of $c \in (a, b)$ such that:
\begin{choices}
    \choice $f(c) = 0$
    \choice $f'(c) = 0$
    \correctchoice $f'(c) = \dfrac{f(b) - f(a)}{b - a}$
    \choice $f(c) = \dfrac{f(a) + f(b)}{2}$
\end{choices}

\question[1] For the function $f(x) = x^{1/3}$, which is true at $x = 0$?
\begin{choices}
    \choice $f$ has a local minimum at $x = 0$.
    \correctchoice $f$ has an inflection point at $x = 0$.
    \choice $f$ has a local maximum at $x = 0$.
    \choice $f$ is not defined at $x = 0$.
\end{choices}

\question[1] $f(x) = x^2 - 4x + 3$ has its minimum value at $x = $
\begin{choices}
    \choice $-2$
    \choice $0$
    \correctchoice $2$
    \choice $4$
\end{choices}

\question[1] The graph of $f$ has $f'(x) > 0$ for $x < 2$ and $f'(x) < 0$ for $x > 2$.
At $x = 2$, $f$ has a:
\begin{choices}
    \choice Local minimum
    \choice Inflection point
    \correctchoice Local maximum
    \choice Vertical asymptote
\end{choices}

\question[1] Which of the following is NOT necessarily true at a critical point $x = c$?
\begin{choices}
    \choice $f'(c) = 0$ or $f'(c)$ is undefined
    \correctchoice $f$ has a local extremum at $x = c$
    \choice $x = c$ could be an inflection point
    \choice Further testing is required to classify $x = c$
\end{choices}

\question[1] At an inflection point of $f$:
\begin{choices}
    \choice $f$ must have a local extremum
    \choice $f'$ must equal zero
    \correctchoice $f''$ changes sign
    \choice $f$ is discontinuous
\end{choices}

\question[1] For $f(x) = x^4$, the point $x = 0$ is a:
\begin{choices}
    \choice Local maximum
    \choice Inflection point
    \correctchoice Local (and absolute) minimum
    \choice The second derivative test is inconclusive, so nothing can be said
\end{choices}


\section*{Part B --- Short Answer (35 marks)}

\question Let $f(x) = x^3 - 3x^2 - 9x + 5$.
\begin{parts}
    \part[2] Find all critical points of $f$.
    \part[2] Use the first derivative test to classify each critical point as a local max,
    local min, or neither.
    \part[2] Find all inflection points.
    \part[2] State the intervals on which $f$ is increasing, decreasing, concave up,
    and concave down.
    \part[2] Sketch the curve, labelling all critical points, inflection points, and the
    $y$-intercept.
\end{parts}
\begin{solution}[11cm]
\textbf{(a)} $f'(x) = 3x^2 - 6x - 9 = 3(x-3)(x+1) = 0 \implies \mathbf{x = -1,\ x = 3}$

$f(-1) = -1 - 3 + 9 + 5 = 10$;\quad $f(3) = 27 - 27 - 27 + 5 = -22$.
Critical points: $(-1,\,10)$ and $(3,\,-22)$.

\textbf{(b)} Sign chart of $f'$:
\begin{itemize}
    \item $x < -1$: $f' > 0$ (increasing);\quad $x \in (-1,3)$: $f' < 0$ (decreasing);\quad $x > 3$: $f' > 0$ (increasing).
    \item $f'$ changes $+\to-$ at $x=-1$: \textbf{local max} at $(-1,\,10)$.
    \item $f'$ changes $-\to+$ at $x=3$: \textbf{local min} at $(3,\,-22)$.
\end{itemize}

\textbf{(c)} $f''(x) = 6x - 6 = 0 \implies x = 1$;\quad $f(1) = 1-3-9+5 = -6$.
$f''$ changes sign at $x=1$ (neg to pos). \textbf{Inflection point at $(1,\,-6)$}.

\textbf{(d)} Increasing: $(-\infty,-1)\cup(3,\infty)$;\quad Decreasing: $(-1,3)$.
Concave down: $(-\infty,1)$;\quad Concave up: $(1,\infty)$.

\textbf{(e)} Sketch: local max $(-1,10)$, inflection $(1,-6)$, local min $(3,-22)$, $y$-intercept $(0,5)$.
\end{solution}

\question An open-top box is to be made by cutting equal squares of side $x$ from the
corners of a $24\;\text{cm} \times 24\;\text{cm}$ piece of cardboard and folding up the sides.
\begin{parts}
    \part[2] Write a function $V(x)$ for the volume of the box. State its domain.
    \part[3] Find the value of $x$ that maximizes the volume. Justify using a derivative test.
    \part[2] State the maximum volume and the dimensions of the resulting box.
\end{parts}
\begin{solution}[9cm]
\textbf{(a)} $V(x) = x(24 - 2x)^2$. Domain: $0 < x < 12$.

\textbf{(b)}
\[
V'(x) = (24-2x)^2 + x \cdot 2(24-2x)(-2) = (24-2x)\bigl[(24-2x) - 4x\bigr]
= (24-2x)(24-6x)
\]
$V'(x) = 0$: $x = 12$ (endpoint) or $x = 4$.

$V''(x) = -2(24-6x) + (24-2x)(-6) = -12(24-2x) + (24-6x)(-2)$... or simply check signs:
$V' > 0$ for $x < 4$, $V' < 0$ for $4 < x < 12$ $\implies$ \textbf{local max at $x = 4$}.

\textbf{(c)} $V(4) = 4(24-8)^2 = 4(256) = \mathbf{1024 \text{ cm}^3}$.
Dimensions: $\mathbf{16 \text{ cm} \times 16 \text{ cm} \times 4 \text{ cm}}$.
\end{solution}

\question A company manufactures cylindrical cans with volume $500\;\text{cm}^3$.
The cost per unit area of the circular top and bottom is twice that of the curved side.
Let $k$ be the cost per cm$^2$ of the side.
\begin{parts}
    \part[2] Write the total cost $C$ as a function of $r$ only.
    (Hint: $V = \pi r^2 h$; side area $= 2\pi r h$; top/bottom area $= 2\pi r^2$.)
    \part[3] Find the value of $r$ that minimizes cost. Verify it is a minimum.
    \part[2] Find the corresponding height $h$ and the ratio $h\!:\!r$.
    \part[1] Interpret the ratio $h\!:\!r$ geometrically.
\end{parts}
\begin{solution}[9cm]
\textbf{(a)} $h = \dfrac{500}{\pi r^2}$.\quad
Side cost: $k \cdot 2\pi r h$;\quad Top/bottom cost: $2k \cdot 2\pi r^2$.
\[
C(r) = k \cdot 2\pi r \cdot \frac{500}{\pi r^2} + 4k\pi r^2 = \frac{1000k}{r} + 4k\pi r^2
\]

\textbf{(b)}
\[
C'(r) = -\frac{1000k}{r^2} + 8k\pi r = 0 \implies 8\pi r^3 = 1000
\implies r^3 = \frac{125}{\pi} \implies r = \frac{5}{\pi^{1/3}} \approx 3.41 \text{ cm}
\]
$C''(r) = \dfrac{2000k}{r^3} + 8k\pi > 0$ for all $r > 0$ $\implies$ \textbf{minimum confirmed}.

\textbf{(c)} $h = \dfrac{500}{\pi r^2}$. From $r^3 = 125/\pi$: $\pi r^3 = 125$, so
$\pi r^2 = 125/r$.
\[
h = \frac{500}{\pi r^2} = \frac{500r}{125} = 4r \approx 13.6 \text{ cm}
\]
$h\!:\!r = 4\!:\!1$.

\textbf{(d)} The optimal can has a height equal to four times its radius (i.e.\ $h = 2d$
where $d$ is the diameter). This reflects the higher cost of the circular faces.
\end{solution}

\question \textbf{Mean Value Theorem.}
\begin{parts}
    \part[2] State the Mean Value Theorem precisely, including all hypotheses.
    \part[4] For $f(x) = x^3 - x$ on $[0,\,2]$, find all values of $c$ in $(0,2)$
    guaranteed by the MVT.
    \part[2] Explain geometrically what the MVT guarantees in terms of secant and tangent lines.
\end{parts}
\begin{solution}[8cm]
\textbf{(a)} If $f$ is continuous on $[a,b]$ and differentiable on $(a,b)$, then there
exists at least one $c \in (a,b)$ such that
\[f'(c) = \frac{f(b) - f(a)}{b - a}.\]

\textbf{(b)} $f(0) = 0$,\quad $f(2) = 8 - 2 = 6$.\quad Average slope $= \dfrac{6}{2} = 3$.

Set $f'(c) = 3$: $\quad 3c^2 - 1 = 3 \implies c^2 = \dfrac{4}{3} \implies c = \dfrac{2}{\sqrt{3}}
= \dfrac{2\sqrt{3}}{3} \approx 1.15$.

Since $\dfrac{2\sqrt{3}}{3} \in (0,2)$, this is the required value: $\boxed{c = \dfrac{2\sqrt{3}}{3}}$.

\textbf{(c)} The MVT guarantees at least one point on the curve where the \textbf{tangent line}
is parallel to (has the same slope as) the \textbf{secant line} joining $(a, f(a))$ to
$(b, f(b))$.
\end{solution}

\end{questions}
\end{document}
